\begin{abstract}
Hallucinations pose critical risks for large language model (LLM)-based agents, often manifesting as hallucinative actions resulting from fabricated or misinterpreted information within the cognitive context. While recent studies have exposed such failures, existing evaluations remain fragmented and lack a principled testbed. In this paper, we present \textbf{MIRAGE-Bench}---\textbf{M}easuring \textbf{I}llusions in \textbf{R}isky \textbf{AGE}nt settings---the first unified benchmark for eliciting and evaluating hallucinations in interactive LLM-agent scenarios. We begin by introducing a three-part taxonomy to address agentic hallucinations: actions that are unfaithful to (i) task instructions, (ii) execution history, or (iii) environment observations. To analyze, we first elicit such failures by performing a systematic audit of existing agent benchmarks, then synthesize test cases using a snapshot strategy that isolates decision points in deterministic and reproducible manners. To evaluate hallucination behaviors, we adopt a fine-grained-level LLM-as-a-Judge paradigm with tailored risk-aware prompts, enabling scalable, high-fidelity assessment of agent actions without enumerating full action spaces. \textsc{MIRAGE-Bench} provides actionable insights on failure modes of LLM agents and lays the groundwork for principled progress in mitigating hallucinations in interactive environments. Our code and dataset are available at \texttt{https://github.com/sunblaze-ucb/mirage-bench.git}.
\end{abstract}